% Part: lambda-calculus
% Chapter: syntax
% Section: eta

\documentclass[../../../include/open-logic-section]{subfiles}

\begin{document}

\olfileid[fr]{lam}{syn}{eta}
\olsection{\texorpdfstring{$\eta$}{η}-conversion}

Il existe une autre relation sur les
lambda-termes. Dans la \olref[fv]{sec}, nous
avons considéré $\lambd[x][(fx)]$, qui
reçoit un argument et lui applique~$f$.
Si $x$ n'est pas libre dans~$f$, ce terme
représente donc la même fonction que~$f$ :
$\lambd[x][(fx)]N$ et $fN$ se réduisent
tous deux à~$fN$. La réduction~$\eta$
et son inverse, l'extension~$\eta$,
formalisent cette idée.

\begin{defn}[Contraction~$\eta$, $\eredone$]
\ollabel{defn:beredone}
La \emph{contraction~$\eta$} ($\eredone$)
est la plus petite relation compatible sur
les termes qui satisfait la condition :
\[
  \lambd[x][M x] \eredone M
  \quad\text{si }x \notin \FV{M}.
\]
\end{defn}

\begin{defn}[Réduction~$\beta\eta$, $\bered$]
\ollabel{defn:bered}
La \emph{réduction~$\beta\eta$} ($\bered$)
est la plus petite relation réflexive et
transitive sur les termes qui contient
$\bredone$ et $\eredone$. Outre la
réflexivité et la transitivité, elle
comprend les deux règles suivantes :
\begin{enumerate}
\item Si $M \bredone N$, alors $M \bered N$.
  \ollabel{defn:bered3}
\item Si $M \eredone N$, alors $M \bered N$.
  \ollabel{defn:bered4}
\end{enumerate}
\end{defn}

\begin{defn}
Nous étendons la relation d'équivalence
$\equal$ par la règle de conversion~$\eta$ :
\[
  \lambd[x][M x] \equal M
  \quad\text{si }x \notin \FV{M}.
\]
La relation étendue se note $\equal[\eta]$.
\end{defn}

La condition $x\notin\FV{M}$ empêche la contraction
d'une abstraction dont le paramètre est libre dans $M$.
Dans $\lambd[x][x x]$, les deux occurrences de $x$
dans le corps sont liées par la même abstraction.
Contracter ce terme en $x$ rendrait une occurrence
de $x$ libre.

L'$\eta$-équivalence est importante par son
lien avec l'extensionnalité des lambda-termes.

\begin{defn}[Extensionnalité]
Nous étendons la relation d'équivalence
$\equal$ par la règle (\ext) :
\begin{center}
Si $Mx \equal Nx$, alors $M \equal N$,
à condition que $x \notin \FV{MN}$.
\end{center}
La relation étendue se note $\equal[\ext]$.
\end{defn}

En termes informels, deux termes vus comme
des fonctions sont égaux lorsqu'ils se
comportent de la même manière pour un même
argument frais.

Nous allons montrer que la règle~$\eta$
exprime exactement cette extensionnalité,
sans ajouter d'autres égalités.

\begin{thm}
$M \equal[\ext] N$ si et seulement si
$M \equal[\eta] N$.
\end{thm}

\begin{proof}
Montrons d'abord, par induction sur la
dérivation, que $\equal[\eta]$ est stable
par la règle d'extensionnalité. Supposons
une application de (\ext) dont la prémisse
est $Mx \equal[\eta] Nx$ par hypothèse
d'induction, avec $x \notin \FV{MN}$.
La compatibilité de $\equal[\eta]$ donne
$\lambd[x][Mx] \equal[\eta]
\lambd[x][Nx]$. Comme $x$ n'est libre
ni dans $M$ ni dans $N$, une contraction~$\eta$
de chaque côté donne $M \equal[\eta] N$.
Ainsi, $\equal[\eta]$ contient
$\equal[\ext]$.

Réciproquement, montrons que la règle~$\eta$
est déjà admise par $\equal[\ext]$.
Si $x \notin \FV{M}$, la contraction~$\beta$
donne $(\lambd[x][Mx])x \equal[\ext] Mx$.
De plus, $x$ n'est libre ni dans
$\lambd[x][Mx]$ ni dans~$M$. La règle (\ext)
donne donc $\lambd[x][Mx] \equal[\ext] M$.
La relation $\equal[\ext]$ contient ainsi
la règle~$\eta$, d'où l'égalité annoncée
des deux relations.
\end{proof}

\end{document}
